\documentclass[11pt,a4paper]{article}
\usepackage[T1]{fontenc}
\usepackage{lmodern,amsmath,amssymb,amsthm,mathtools,booktabs,microtype}
\usepackage[margin=27mm]{geometry}
\usepackage[colorlinks=true,linkcolor=blue,urlcolor=blue,citecolor=blue]{hyperref}
\newtheorem{theorem}{Theorem}[section]
\newtheorem{lemma}[theorem]{Lemma}
\newtheorem{proposition}[theorem]{Proposition}
\newtheorem{corollary}[theorem]{Corollary}
\newtheorem{definition}[theorem]{Definition}
\newcommand{\Q}{\mathbb Q}\newcommand{\Z}{\mathbb Z}
\newcommand{\A}{\mathbb A}
\DeclareMathOperator{\Norm}{N}
\DeclareMathOperator{\Frac}{Frac}
\DeclareMathOperator{\Gal}{Gal}
\DeclareMathOperator{\disc}{disc}
\setlength{\emergencystretch}{2em}
\title{Optimal polynomial square tests in quadratic number fields}
\author{}\date{8 November 2026}
\begin{document}\maketitle
\begin{abstract}
For every quadratic extension $K/F$ of number fields, square membership
in $K$, with two input components in $F$, is expressible by one polynomial
of total degree six and one auxiliary $F$-variable. Zero is included and
no nonzero condition is needed. We prove that total degree six and degree
four in the auxiliary variable are both minimal for this form of
definition. The lower bound uses Hilbert irreducibility, the intermediate
fields of the generic squaring map, and a valuation formula for leading
coefficients. Over $\Q(\sqrt2)$ we give an integer polynomial with ten
monomials and coefficient height $7$. A degree-ten alternative has
coefficient height $8$; allowing two auxiliary variables gives a quartic
of height $8$, or two quadratic equations. All results in this paper are
unconditional. They concern square membership, rather than integrality.
\end{abstract}

\section{Inputs, quantified variables, and the main results}

Fix a number field $F$, an element $d\in F^\times\setminus F^{\times2}$,
and $K=F(\rho)$, where $\rho^2=d$. The two input coordinates $U,V$
describe the element $U+V\rho$ of $K$. Write
\[
 \mathcal S_d(F)=\{(U,V)\in F^2:U+V\rho\text{ is a square in }K\}.
\]
Throughout, zero is a square. The scalar $d$ is fixed, not a variable.

\begin{definition}
A polynomial square test with one auxiliary variable is a nonzero
$P\in F[U,V,T]$ such that, for every $(U,V)\in F^2$,
\begin{equation}\label{eq:predicate}
 (U,V)\in\mathcal S_d(F)
 \quad\Longleftrightarrow\quad
 \exists T\in F\quad P(U,V,T)=0.
\end{equation}
\end{definition}

Thus the polynomial has three total coordinates: two inputs and one
quantified variable. All three coordinates range over the base field
$F$. This differs from a polynomial with one input and two auxiliary
variables defining $\Z$ over $\Q$.

\begin{theorem}\label{thm:main}
For every quadratic extension of number fields as above, there is a
square test of total degree six, of degree four in $T$, and with a
nonzero constant leading coefficient in $T$. Every square test
\eqref{eq:predicate}, whether or not its leading coefficient is constant,
satisfies
\[
                         \deg P\ge6,\qquad \deg_TP\ge4.
\]
No polynomial equation without auxiliary variables defines
$\mathcal S_d(F)$.
\end{theorem}

The construction is proved in Sections~\ref{sec:explicit} and
\ref{sec:general}; the lower bounds are proved in
Section~\ref{sec:optimality}. Over $\Q(\sqrt2)$ the following alternatives
make the degree and variable tradeoff explicit:
\begin{center}
\begin{tabular}{lrrrr}
\toprule
Presentation & Auxiliary variables & Total degree & Monomials & Height\\
\midrule
One sextic equation & 1 & 6 & 10 & 7\\
One degree-ten equation & 1 & 10 & 12 & 8\\
One quartic equation & 2 & 4 & 8 & 8\\
Two quadratic equations & 2 & 2 & --- & 2\\
\bottomrule
\end{tabular}
\end{center}
For a system, the displayed degree and coefficient height are the
maxima over its equations. Coefficient height means the largest absolute
integer coefficient. These are square tests; none of the displayed
equations is asserted to define a ring of integers.

\section{An explicit sextic over the rationals}\label{sec:explicit}

Put
\begin{equation}\label{eq:AB}
 A=3U^3-2U^2V-3UV^2,\qquad
 B=-2U^3+U^2V+2UV^2+V^3,
\end{equation}
and define
\begin{equation}\label{eq:G}
 \boxed{G(U,V,T)=2T^4-2A(U,V)T^2+B(U,V)^2.}
\end{equation}

\begin{lemma}\label{lem:B}
For rational $U,V$, one has $B(U,V)=0$ if and only if $U=V=0$.
\end{lemma}
\begin{proof}
Scale a hypothetical nonzero rational zero to relatively prime
integers $U,V$ and reduce modulo $3$. The values of $B(t,1)$
at $t=0,1,2$ are $1,2,2$, respectively, and $B(1,0)=1$
modulo $3$. Thus $B$ has no zero in $\mathbf P^1(\mathbf F_3)$,
a contradiction.
\end{proof}

\begin{theorem}\label{thm:explicit}
For every $U,V\in\Q$,
\[
 U+V\sqrt2\text{ is a square in }\Q(\sqrt2)
 \quad\Longleftrightarrow\quad
 \exists T\in\Q\quad G(U,V,T)=0.
\]
The polynomial has total degree six, degree four in $T$, ten monomials,
integer content one, and coefficient height $7$.
\end{theorem}
\begin{proof}
Define the multiplier
\[
             m=(U+V)-U\sqrt2.
\]
Direct multiplication gives
\begin{align}
 \Norm_{K/\Q}(m)&=-U^2+2UV+V^2,\label{eq:multnorm}\\
 m^2(U+V\sqrt2)&=A+B\sqrt2.\label{eq:multsquare}
\end{align}
If $U+V\sqrt2=(u+v\sqrt2)^2$, then
\begin{equation}\label{eq:cubic}
 m(u+v\sqrt2)=(u^3+2uv^2-4v^3)+(-u^3+u^2v+2v^3)\sqrt2.
\end{equation}
For a nonzero input, Lemma~\ref{lem:B} gives $B\ne0$.
The product of the two components on the right of~\eqref{eq:cubic}
is $B/2$, so neither component is zero. Taking
$T=u^3+2uv^2-4v^3$ and comparing the square's components gives
\[
               A=T^2+\frac{B^2}{2T^2},
\]
which is exactly $G=0$. The zero square is represented by $U=V=T=0$.

Conversely, suppose $G(U,V,T)=0$. If $T=0$, Lemma~\ref{lem:B}
gives $U=V=0$. Otherwise
\[
 \left(T+\frac{B}{2T}\sqrt2\right)^2=A+B\sqrt2.
\]
For $(U,V)\ne(0,0)$ the norm in~\eqref{eq:multnorm} is nonzero:
$(U+V)^2=2U^2$ would force $U=V=0$.
Division by $m$ therefore gives
the explicit square root
\begin{equation}\label{eq:recover}
 z=\frac{(U+V)+U\sqrt2}{-U^2+2UV+V^2}
                         \left(T+\frac{B}{2T}\sqrt2\right).
\end{equation}
There is no missing exceptional input: $T=0$ was handled separately,
and every nonzero input has a nonzero multiplier.

The full expansion is
\begin{align*}
 G={}&4U^6-4U^5V-7U^4V^2+6U^2V^4+4UV^5+V^6\\
    &-6U^3T^2+4U^2VT^2+6UV^2T^2+2T^4.
\end{align*}
It proves the stated degree, monomial count, content, and height.
\end{proof}

The polynomial is quartic in its auxiliary variable, but it is
\emph{sextic in total degree}. For example, $U^6$ has degree six.
No matrix entry, component, or multiplier in the displayed formula
is an additional quantified coordinate; they are explicit polynomials
in the two given inputs.

\begin{proposition}\label{prop:separable}
For each nonzero rational input $(U,V)$, every algebraic root of
$G(U,V,T)$ is nonzero and simple. Its discriminant in $T$ is
\begin{equation}\label{eq:disc}
 \disc_TG=512B^2(-U^2+2UV+V^2)^4(U^2-2V^2)^2.
\end{equation}
\end{proposition}
\begin{proof}
Taking norms in~\eqref{eq:multsquare} gives
\[
        A^2-2B^2=(-U^2+2UV+V^2)^2(U^2-2V^2).
\]
The discriminant of $aT^4+bT^2+c$ is
$16ac(b^2-4ac)^2$, obtained from its derivative and resultant.
Applying this identity gives~\eqref{eq:disc}. Each factor is nonzero:
use Lemma~\ref{lem:B} and the nonsquareness of $2$. The constant
term $B^2$ is also nonzero.
\end{proof}

\subsection{Coefficient height within the integral multiplier family}

The following coefficient lower bound has a narrower scope than
the degree lower bound in Theorem~\ref{thm:main}.

\begin{proposition}\label{prop:height7}
Let $a,b,c,e$ be integers and put
\[
 p=aU+bV,\qquad q=cU+eV,\qquad
 A_m=(p^2+2q^2)U+4pqV,\quad
 B_m=(p^2+2q^2)V+2pqU.
\]
Divide $2T^4-2A_mT^2+B_m^2$ by its positive integer content.
If the resulting polynomial is a complete square test, its
coefficient height is at least seven. Equality is attained by
\eqref{eq:G}.
\end{proposition}
\begin{proof}
The four coefficients of $B_m$, in descending $U$-degree, are
\[
 2ac,\quad a^2+2c^2+2ae+2bc,\quad
 2ab+4ce+2be,\quad b^2+2e^2.
\]
Soundness requires $ac\ne0$: otherwise every $(U,0,0)$ would
be a zero, although some rational $U$ are not squares in $K$.
It also requires $(b,e)\ne(0,0)$, since otherwise every
$(0,V,0)$ would be a zero. For $V\ne0$, the norm of $V\sqrt2$
is negative, so that element cannot be a square in $K$.

The content divides two, the coefficient of $T^4$. Modulo two,
$B_m\equiv a^2U^2V+b^2V^3$; hence the content is two exactly
when $a,b$ are both even. In that case the primitive $U^6$
coefficient is $2a^2c^2\ge8$.

Suppose instead that the content is one and the height is at
most six. The $U^3T^2$ coefficient $-2(a^2+2c^2)$ forces
$a,c\in\{1,-1\}$. The $V^6$ coefficient $(b^2+2e^2)^2$
then forces
\[
 (b,e)\in\{(1,0),(-1,0),(0,1),(0,-1)\}.
\]
Write $B_j=[U^{3-j}V^j]B_m$. Then $B_0=2ac$ and
$B_1=3+2ae+2bc\in\{1,5\}$. The $U^5V$ coefficient is
$2B_0B_1$, so the proposed height bound requires $B_1=1$.
If $e=0$, this implies $bc=-1$ and $B_2=-2ac$; the
$U^4V^2$ coefficient is therefore $1-8=-7$.
If $b=0$, one has $ae=-1$ and $B_2=-4ac$, giving $1-16=-15$.
Both conclusions contradict height at most six.
Theorem~\ref{thm:explicit} proves attainment.
\end{proof}

This proposition does not assert coefficient optimality among
all square tests or after arbitrary rational changes of the
auxiliary coordinate.

\section{A degree-ten alternative}\label{sec:ten}

The alternative polynomial
\begin{equation}\label{eq:P10}
 P_{10}(U,V,T)
   =(2-2UT^2+V^2T^4)\bigl(V^2+(U-2T^2)^2\bigr)
\end{equation}
has total degree ten, degree eight in $T$, twelve monomials, and
coefficient height eight. It too is a complete square test, including zero.

To prove this, suppose first that its first factor vanishes. Then $T\ne0$
and
\[
       \left(\frac1T+\frac{VT}{2}\sqrt2\right)^2=U+V\sqrt2.
\]
Conversely a square root $u+v\sqrt2$ with $u\ne0$ gives this factor
zero with $T=1/u$. The second factor vanishes precisely when
$V=0$ and $U=2T^2$, because a sum of rational squares vanishes only
when both squares vanish. These are exactly the inputs with a square
root $T\sqrt2$, including zero. A product is zero if and only if one
factor is zero, so the two cases use the same one quantified variable.

The highest-degree monomial is $4V^2T^8$, of total degree ten.
The value eight is the largest coefficient after complete expansion;
it is not an estimate for the coefficients in any integrality construction.

Over a field $F$ of characteristic different from two, the same
argument gives
\begin{equation}\label{eq:general10}
 (4-4UT^2+dV^2T^4)\bigl(V^2-c(U-dT^2)^2\bigr)=0
\end{equation}
whenever $d$ and $c$ are nonzero nonsquares in $F$. One may take $c=d$.
The second factor then forces $V=0$ and $U=dT^2$.

\section{A sextic over every quadratic number field}\label{sec:general}

Choose $\kappa\in F^\times$ that is not a cube. Such an element exists
in every number field: choose a nonzero prime ideal $\mathfrak p$ and
$\kappa\in\mathfrak p\setminus\mathfrak p^2$.
Define
\begin{align}
 a&=U+\frac{\kappa}{2d}V,
 &b&=-\frac{\kappa}{d^2}U-\frac12V,\label{eq:generalab}\\
 A_d&=(a^2+db^2)U+2dabV,
 &B_d&=(a^2+db^2)V+2abU.\label{eq:generalAB}
\end{align}
The parameters $d,\kappa$ are fixed coefficients. Thus $a,b$ are linear
and $A_d,B_d$ are homogeneous cubics in the two input variables.

\begin{lemma}\label{lem:generalB}
The cubic $B_d$ vanishes at a pair in $F^2$ only at $(0,0)$.
\end{lemma}
\begin{proof}
Expansion gives
\[
 4d^2B_d=-8\kappa U^3+6\kappa dUV^2+(\kappa^2+d^3)V^3.
\]
If $V=0$, its vanishing implies $U=0$. Otherwise, put $t=2U/V$.
A zero would give an $F$-root of
\begin{equation}\label{eq:irreduciblecubic}
             X^3-3dX-\kappa-\frac{d^3}{\kappa}.
\end{equation}
Let $r^3=\kappa$. The polynomial $X^3-\kappa$ is irreducible over $F$,
so $[F(r):F]=3$. The element $s=r+d/r$ is a root
of~\eqref{eq:irreduciblecubic}. Moreover, $r$ satisfies
$X^2-sX+d=0$ over $F(s)$. Therefore $[F(r):F(s)]$ divides three
and is at most two, so it is one. Thus $s$ has degree three over $F$,
and~\eqref{eq:irreduciblecubic} is irreducible. It has no $F$-root.
\end{proof}

\begin{theorem}\label{thm:general}
For every $U,V\in F$,
\begin{equation}\label{eq:general6}
 (U,V)\in\mathcal S_d(F)
 \quad\Longleftrightarrow\quad
 \exists Y\in F\quad
 Y^4-4A_d(U,V)Y^2+4dB_d(U,V)^2=0.
\end{equation}
This is one polynomial of exact total degree six and degree four in $Y$.
For each nonzero input, all its algebraic roots are nonzero and simple.
\end{theorem}
\begin{proof}
Let $m=a+b\rho$. The identities
\begin{align*}
 m^2(U+V\rho)&=A_d+B_d\rho,\\
 a^2-db^2&=\frac{d^3-\kappa^2}{4d^3}(4U^2-dV^2)
\end{align*}
follow by multiplication. Both factors in the norm are nonzero for
nonzero $(U,V)$: $\kappa^2=d^3$ would make $d=(\kappa/d)^2$ a
square, and $4U^2=dV^2$ would likewise contradict the choice of $d$.

For a square root $z=u+v\rho$, substitution gives
\[
 mz=(u^3-\kappa v^3)
           +\left(dv^3-\frac{\kappa}{d^2}u^3\right)\rho.
\]
The first component is nonzero when $z\ne0$. Set $Y$ equal to twice
that component to obtain the equation in~\eqref{eq:general6}.
Zero is represented by $U=V=Y=0$.

Conversely, a root $Y=0$ forces $B_d=0$ and hence $U=V=0$ by
Lemma~\ref{lem:generalB}. For a nonzero input, $Y\ne0$, and the
displayed equation gives
\[
       \left(\frac Y2+\frac{B_d}{Y}\rho\right)^2=A_d+B_d\rho.
\]
Dividing the element in parentheses by $m$ recovers a square root.
At a permitted root the derivative is $8Y$ times the norm of that
parenthesized element. Its squared norm is
$(a^2-db^2)^2(U^2-dV^2)\ne0$, so the derivative is nonzero.
Finally $B_d$ is a nonzero cubic, so $4dB_d^2$ has exact degree six;
the other terms have smaller total degree and cannot cancel it.
\end{proof}

\section{Why degree six is necessary}\label{sec:optimality}

We give the lower bound for an arbitrary polynomial square test,
without a monicity or irreducibility assumption on that polynomial.
The only external theorem used in this section is Hilbert irreducibility:
over a number field, finitely many irreducible polynomials over a
rational function field can be specialized simultaneously so as to
remain irreducible, away from a prescribed proper closed subset
\cite[Proposition~3.3.5 and Theorem~3.4.1]{Serre}.

\subsection{The generic squaring map and its intermediate fields}

Let $u,v$ now be algebraically independent over $F$ and put
\[
 U=u^2+dv^2,\quad V=2uv,\quad w=u^2-dv^2,
 \quad K_0=F(U,V),\quad L=F(u,v),\quad M=K_0(w).
\]
These function fields are distinct from the fixed quadratic number field
$K=F(\rho)$.

\begin{lemma}\label{lem:fields}
The extension $L/K_0$ has degree four. Its unique proper intermediate
field is $M$, with $[M:K_0]=2$. The integral closure of $F[U,V]$ in $L$
is $F[u,v]$.
\end{lemma}
\begin{proof}
Over $F(\rho)$ set $x=u+\rho v$ and $y=u-\rho v$. Then
$U+\rho V=x^2$ and $U-\rho V=y^2$. The independent sign changes
of $x,y$ give a degree-four extension over $F(\rho)(U,V)$.
The field $N=F(\rho,u,v)$ is a degree-eight normal extension of $K_0$:
besides the sign changes, it has the automorphism interchanging $x,y$
and sending $\rho$ to $-\rho$. Its group is the dihedral group of order
eight. The subgroup fixing $L$ consists of the identity and this
interchange. Its only proper overgroup other than itself has order four,
generated by the interchange and the simultaneous sign change.
Its fixed field is $M$, since $w=xy$ is fixed and
$w^2=U^2-dV^2$. This proves the intermediate-field assertions.

Both $u$ and $v$ are integral over $F[U,V]$, as seen from
\[
 u^4-Uu^2+dV^2/4=0,\qquad
 v^4-(U/d)v^2+V^2/(4d)=0.
\]
Thus $F[u,v]$ is finite over $F[U,V]$. It is integrally closed, being
a polynomial ring, and has fraction field $L$. Any element of $L$
integral over $F[U,V]$ is also integral over $F[u,v]$, hence lies in
$F[u,v]$. This proves the normalization assertion.
\end{proof}

Choose once and for all
\begin{equation}\label{eq:badc}
                    c\in F^\times\setminus(F^{\times2}\cup dF^{\times2}).
\end{equation}
Such a choice is possible by taking an element of valuation one at a
prime ideal where $d$ has valuation zero. Expanding
$(a+b\rho)^2\in F$ shows that $c$ is a square in $K$ if and only if
$c$ or $c/d$ is a square in $F$. Thus our $c$ is not a square in $K$.

\begin{lemma}\label{lem:bad-dense}
The surface $w^2=U^2-dV^2$ has a Zariski-dense set of $F$-points
whose input $U+V\rho$ is not a square in $K$.
\end{lemma}
\begin{proof}
Use the rational parametrization
\[
 a=U+w,\quad V=b,\quad
 U=\frac12\left(a+\frac{db^2}{a}\right),\quad
 w=\frac12\left(a-\frac{db^2}{a}\right),\qquad a\ne0.
\]
Put $a=2cs^2$, where $s\in F^\times$, and let $b\in F$ vary.
If the resulting input were $(p+q\rho)^2$, then
$p^2-dq^2$ would equal $w$ or $-w$. The first possibility gives
$p^2=a/2=cs^2$, impossible. In the second,
$dq^2=a/2=cs^2$, also impossible. The same argument includes $b=0$.
The map $(s,b)\mapsto(a,b)=(2cs^2,b)$ is dominant, and $F^2$ is
Zariski dense. Its images are consequently dense on the surface.
\end{proof}

\begin{lemma}\label{lem:primitive}
If $P$ is a polynomial square test, the polynomial
$P(u^2+dv^2,2uv,T)$ has a root $r\in L$. Every such root generates
$L$ over $K_0$.
\end{lemma}
\begin{proof}
If its irreducible factors over $F(u,v)$ all had degree at least two
in $T$, simultaneous Hilbert irreducibility would give a specialization
$(u,v)\in F^2$ where none had an $F$-root and all required coefficient
denominators were nonzero. The input is a square, contradicting
the completeness of the test. A factor of degree one therefore exists.

Suppose a root $r$ did not generate $L$. By Lemma~\ref{lem:fields},
it would lie in $M$ (which includes the possibility $r\in K_0$).
Regard it as a rational function on the surface in
Lemma~\ref{lem:bad-dense}. Away from its poles, the identity
$P(U,V,r)=0$ specializes to an $F$-root of $P$. The dense set of
nonsquare inputs in that lemma meets this open subset. This contradicts
the soundness of the square test.
\end{proof}

In particular the minimal polynomial of $r$ over $K_0$ has degree four,
so $\deg_TP\ge4$. Choose its primitive polynomial representative
\[
 f(U,V,T)=a(U,V)T^4+\cdots\in F[U,V,T].
\]
Here ``primitive'' means that its coefficients as a polynomial in $T$
have no nonconstant common factor in $F[U,V]$. Gauss's lemma implies
$f\mid P$ in $F[U,V,T]$, and therefore $\deg f\le\deg P$.

\subsection{Poles and leading coefficients}

For an irreducible divisor $D$ of the affine input plane, write $v_D$
for the integer-normalized valuation on $K_0$. Its extensions $w_i$
to $L$ are also normalized to have value group $\Z$; their residue
degrees are denoted by $f_i$.

\begin{lemma}\label{lem:gauss}
With the preceding notation,
\begin{equation}\label{eq:leadingval}
                    v_D(a)=\sum_i f_i\max(0,-w_i(r)).
\end{equation}
\end{lemma}
\begin{proof}
For a polynomial over a valued field, the least valuation of its
coefficients is additive under multiplication. To see this, divide
each factor by one of its coefficients of least valuation and reduce
modulo the maximal ideal; the two nonzero reduced polynomials have
a nonzero product. Apply this fact over a normal closure to the
monic minimal polynomial
$f/a=\prod_{\sigma}(T-\sigma r)$. Its least coefficient valuation is
$\sum_\sigma\min(0,v_D(\sigma r))$.

At a place with ramification index $e_i$ and residue degree $f_i$,
there are $e_if_i$ conjugates, and their extended valuation is
$w_i(r)/e_i$. The sum is therefore
$\sum_i f_i\min(0,w_i(r))$. Since $f$ is primitive, its coefficients
have least $v_D$-valuation zero. Multiplying by $a$ gives
\eqref{eq:leadingval}.
\end{proof}

\begin{lemma}\label{lem:line}
If $D$ is an affine $F$-line in the input plane, every place of $L$
above $D$ has residue degree at least two.
\end{lemma}
\begin{proof}
The map associated with the finite extension
$F[U,V]\subset F[u,v]$ is the squaring map. A residue-degree-one
point above the generic point of $D$ would give $u(t),v(t)\in F(t)$
such that
\[
 u(t)^2+dv(t)^2=U(t),\qquad 2u(t)v(t)=V(t),
\]
where $U(t),V(t)$ are affine linear and not both constant.
In $K(t)$ this would imply
\[
                  (u(t)+\rho v(t))^2=U(t)+\rho V(t).
\]
The right side is a nonconstant linear polynomial: its leading
coefficient cannot vanish because $1,\rho$ are linearly independent
over $F$. It has a simple zero and hence cannot be a rational-function
square. This contradiction proves the assertion.
\end{proof}

\begin{lemma}\label{lem:polynomial-r}
If $\deg f\le5$, then $a$ is constant and $r\in F[u,v]$ has degree
at most two in $u,v$.
\end{lemma}
\begin{proof}
Since $\deg_Tf=4$, its leading coefficient $a$ has degree at most one.
If it were nonconstant, its zero divisor would be an affine line $D$
with $v_D(a)=1$. By~\eqref{eq:leadingval}, some $w_i(r)$ would be
negative. Lemma~\ref{lem:line} would then make the right side of
\eqref{eq:leadingval} at least two, a contradiction. Thus $a$ is
constant. After dividing by it, $f$ is monic, so $r$ is integral over
$F[U,V]$ and belongs to $F[u,v]$ by Lemma~\ref{lem:fields}.

Let $m=\deg_{u,v}r$. The four conjugates of $r$ in the normal
closure are obtained by invertible linear changes of $u,v$ over
$K$, so each has degree $m$. Their product has degree $4m$ in $u,v$.
This product is the norm of $r$, the constant coefficient of its monic
degree-four minimal polynomial. Because $U,V$ are algebraically
independent homogeneous quadratics, a polynomial of degree $e$ in
$U,V$ has degree $2e$ after their substitution. Hence
\[
                     \deg_{U,V}\Norm_{L/K_0}(r)=2m.
\]
The constant coefficient has degree at most five, giving $m\le2$.
\end{proof}

\subsection{The final obstruction}

\begin{theorem}\label{thm:lower}
Every polynomial square test with one auxiliary variable has total
degree at least six.
\end{theorem}
\begin{proof}
Suppose $\deg P\le5$, and choose $r$ and $f$ as above. By
Lemma~\ref{lem:polynomial-r}, write
\[
                  r=q_2(u,v)+a_1u+b_1v+q_0,
\]
where $q_2$ is homogeneous quadratic and $q_0\in F$.
The quadratic and constant terms are fixed by $(u,v)\mapsto(-u,-v)$,
so belong to $M$. Since $r$ generates $L$, we have
$(a_1,b_1)\ne(0,0)$.

Use the nonsquare $c$ from~\eqref{eq:badc} and specialize over an
algebraic closure to
\[
                       u=b_1\sqrt c,\qquad v=-a_1\sqrt c.
\]
Then $U,V\in F$, the linear part of $r$ vanishes, and its quadratic
part has value in $F$. Thus the specialized value $t_0$ of $r$ lies
in $F$. The polynomial identity $f(u^2+dv^2,2uv,r(u,v))=0$ implies
$f(U,V,t_0)=0$, and divisibility $f\mid P$ gives $P(U,V,t_0)=0$.
However,
\[
                   U+V\rho=c(b_1-a_1\rho)^2
\]
is not a square in $K$. The factor $b_1-a_1\rho$ is nonzero, and
$c$ was chosen nonsquare in $K$. This contradicts soundness.
\end{proof}

\begin{proposition}\label{prop:no-aux}
The set $\mathcal S_d(F)$ cannot be defined by polynomial equations
without auxiliary variables, even by a finite union of such systems.
\end{proposition}
\begin{proof}
The squaring map is dominant: after extending scalars to $K$, it is
the map $(x,y)\mapsto(x^2,y^2)$ in the coordinates used in
Lemma~\ref{lem:fields}. Since $F^2$ is Zariski dense, so is the image
$\mathcal S_d(F)$. It is a proper subset of $F^2$, by the nonsquare
choice~\eqref{eq:badc}. A finite union of sets defined by polynomial
equalities is the set of $F$-points of a Zariski-closed set. A closed
set containing a dense subset is the whole plane, which is impossible.
\end{proof}

These results prove all assertions of Theorem~\ref{thm:main}.

\section{Several equations and two auxiliary variables}

\begin{corollary}\label{cor:systems}
Suppose finitely many polynomial equalities, with two input variables
and just one common auxiliary $F$-variable, define $\mathcal S_d(F)$.
Then at least one nonzero equation has total degree at least six.
The conclusion also holds for a finite disjunction of such systems.
\end{corollary}
\begin{proof}
For a system, form the greatest common divisor of its nonzero
polynomials in $K_0[T]$. If that divisor were one, a Bezout identity
would exclude all common roots outside a proper closed subset of the
input plane, contrary to square-input completeness. After pulling
back by the squaring map and avoiding coefficient denominators, common
roots are roots of this gcd. Hilbert irreducibility therefore supplies
a root $r\in L$, as in Lemma~\ref{lem:primitive}. Its primitive minimal
polynomial divides every nonzero equation of the system. All $F$-roots
of that factor consequently satisfy the system, so the soundness and
degree arguments of Section~\ref{sec:optimality} apply unchanged.

For a finite disjunction, apply the same argument to the gcd of each
branch. Branches with gcd one can occur only on a proper closed set.
If every remaining pullback factor had degree at least two, one
simultaneous Hilbert specialization would avoid rational roots of all
branches. Thus one branch has a root in $L$; its minimal polynomial
divides all equations in that branch, and the same contradiction applies.
\end{proof}

The assertion concerns polynomial \emph{equalities}. Nonzero conditions
change the situation: the quartic $4T^4-4UT^2+dV^2=0$, together
with $T\ne0$, covers the squares with a square root whose first
component is nonzero. The missing case is a second branch
$V=0$, $U=dT^2$. Replacing a nonzero condition by an equality
usually requires an additional variable, which must be counted.

With two auxiliary variables there is a simpler complete presentation:
\begin{equation}\label{eq:twoquad}
                   U=u^2+dv^2,\qquad V=2uv.
\end{equation}
Over $\Q$, these two equations combine without further variables into
\begin{equation}\label{eq:quartic}
             (U-u^2-2v^2)^2+(V-2uv)^2=0.
\end{equation}
This is a quartic in four total coordinates, with coefficient height
eight and eight monomials. Over a general number field $F$, choose
a nonsquare $c_0$ and combine the residuals $R_1,R_2$ as
$R_1^2-c_0R_2^2=0$; this is again a quartic with the same four
coordinates. Thus allowing more variables can lower the degree,
while merely adding equations does not lower the sharp bound with
one auxiliary variable.

\section{Relation to the elliptic and integrality constructions}

For $D=(3-\sqrt2)^2$, direct elliptic charts use
\[
       x=h^2,\qquad j^2=h^4-D,\qquad y=hj,
       \qquad h=h_0+h_1\sqrt2.
\]
The two input components of the square-root equation are
\[
 U=h_0^4+12h_0^2h_1^2+4h_1^4-11,\qquad
 V=4h_0h_1(h_0^2+2h_1^2)+6.
\]
They cannot both vanish. Otherwise $h^2=\pm(3-\sqrt2)$;
the negative sign is excluded by a real embedding and the positive
sign by the nonsquare norm $7$. Proposition~\ref{prop:separable}
and the general construction in Theorem~\ref{thm:general}
therefore give nonzero, simple finite roots at every rational
specialization of $h_0,h_1$. This justifies treating the corresponding
root, multiplier-norm, and derivative conditions as automatic when
their only purpose is to guard this chart.

The lower bound in Theorem~\ref{thm:main} concerns \emph{all} input
pairs $(U,V)\in F^2$. It does not automatically give the same bound
for a restricted family of radicands, or for an entire integrality
construction in which different algebraic conditions interact.

In particular, the equations in this paper test square membership,
not membership in $\Z$ or in a ring of integers. For example, $1/4$
is a square in $\Q(\sqrt2)$ but is not an algebraic integer, whereas
$3$ is an algebraic integer and is not a square in that field. The
last assertion follows by expanding $(u+v\sqrt2)^2=3$: either
$v=0$ and $u^2=3$, or $u=0$ and $v^2=3/2$, both impossible over $\Q$.

The ring of integers of $\Q(\sqrt2)$ is $\Z[\sqrt2]$. For completeness,
if $a+b\sqrt2$ is integral, then $2a=m\in\Z$ and
$a^2-2b^2=n\in\Z$. From $8b^2=m^2-4n$ a reduced denominator
of $b$ has square dividing $8$, so $b=k/2$ for an integer $k$.
Then $m^2-2k^2\equiv0\pmod4$ forces both $m$ and $k$ even.
Thus $a,b\in\Z$; the converse is immediate. Defining this ring
existentially inside its fraction field is a different problem.

All auxiliary-variable lower bounds here use quantification over $F$.
With an auxiliary variable $z$ ranging over $K$, square membership
already has the simple equation $z^2=U+V\rho$ over $K$.
Our formulas express the same condition using only one $F$-coordinate,
instead of the two $F$-components of $z$. They do not assert a
one-coordinate transfer of arbitrary existential formulas, or an
integrality definition over $K$.

\section{Exact computational data}

The accompanying standard-library Python program
\texttt{check\_square\_predicates.py} reproduces the complete
coefficient tables for~\eqref{eq:G}, \eqref{eq:P10}, and
\eqref{eq:quartic}. It independently checks their expanded formulas,
degrees, integer content, and coefficient heights. It checks the
general multiplier and norm identities by exact rational polynomial
arithmetic, and computes the seven-by-seven Sylvester determinant
for~\eqref{eq:disc}, rather than relying only on a quoted discriminant
formula. A finite permutation-group enumeration verifies the
overgroup calculation used in Lemma~\ref{lem:fields}.

The generated files are
\begin{quote}\small
\nolinkurl{square-predicate-polynomials.json}\\
\nolinkurl{square-predicate-verification.json}.
\end{quote}
Their calculations are
finite exact certificates for the listed algebraic assertions.
Hilbert irreducibility and the uniform degree lower bound are used
and proved, respectively, in the text; numerical examples are not
substituted for either argument.

\begin{thebibliography}{9}
\bibitem{Serre}
J.-P.~Serre, \emph{Topics in Galois Theory}, second edition,
Research Notes in Mathematics, vol.~1, A K Peters, Wellesley, MA,
2007. Notes by H.~Darmon. Proposition~3.3.5 and Theorem~3.4.1,
pp.~24--25.
\url{https://indico.math.cnrs.fr/event/11410/attachments/4760/7351/SerreTopicsGaloisTheory.pdf}.
\end{thebibliography}
\end{document}
